\section{Conclusion}
\label{sec:conc}
A zero-parameter, interpretable, accelerator-free method---used both as a competitor and as a diagnostic
instrument---is the strongest public-feature approach to large-scale bundle completion on $5$ of $8$ benchmarks: it
beats CLHE and RaMen on POG (${\sim}1.7\times$ R@20), on Spotify ($0.1562$ vs.\ $0.1091$), and on Spotify\_sparse
($0.1653$ vs.\ $0.0969$), wins on Clothing and Dunnhumby, ties the $9.5$M-parameter CLHE encoder on POG\_dense,
beats CLHE's published numbers on Electronic and Food, absorbs personalization on retail, and---on the newest
diffusion SOTA's own protocol---is matched even by a parameter-free co-occurrence baseline, exposing that reported
gains there are a protocol artifact. Its only shortfall is against RaMen on the two dense e-commerce sets, and---from
RaMen's own ablation---we trace that gap to its learned multi-hop collaborative encoder (whose removal drops RaMen
below us), not to its product images (whose removal barely moves it, and whose $39\%$-coverage public recovery does
not help us either). We report the negatives in the open---inert higher-order operators, a rejected homoscedastic
variance model, and an unhelpful image signal---and verify that validation
weight-selection does not
overfit on any dataset. The message is not ``cheap always wins,'' but the more useful one: \emph{a method that fits
the regime is competitive with heavy GPU/diffusion/LLM machinery wherever the comparison is actually checkable---know
your regime, and spend heavy compute only where it provably buys something.} A natural next step is to turn the
closed-form score into a \emph{calibrated} decision model---emitting thresholdable membership probabilities via a
proper-scoring-rule fit---so a recommender can act automatically above a confidence level and abstain below it.
